\phantomsection
\subsection*{M2D. Stone Fruit Mead}
\addcontentsline{toc}{subsection}{M2D. Stone Fruit Mead}

\textit{Stone Fruit Mead é uma categoria de inscrição para meloméis elaborados com uma ou mais frutas de caroço, como cereja, ameixa, pêssego, damasco e manga. Frutas de caroço são frutas carnudas com um único caroço grande.}

\textbf{Impressão Geral}: Assemelha-se a um M1 Traditional Mead, porém com um caráter perceptível de fruta. Uma ampla gama de resultados é possível dependendo das escolhas do hidromel base e dos ingredientes adicionais, mas o produto final deve ser equilibrado, agradável e prazeroso. O caráter da fruta não deve dominar completamente o hidromel.

\textbf{Aroma}: Os aromas da fruta de caroço podem ser percebidos como de um suco doce e fresco até um vinho de fruta envelhecido ou do Porto, dependendo da força alcoólica e do dulçor. Se variedades de fruta forem declaradas, o caráter varietal deve estar perceptível. Se múltiplos tipos de fruta forem declarados, eles não precisam estar em proporção igual, tendo em vista que algumas frutas são mais fortes e distintivas que outras. O caráter do mel pode variar de sutil a intenso, dependendo da força e dulçor, podendo expressar o néctar floral típico de eventuais variedades de mel declaradas. O buquê de fermentação deve ser equilibrado. Hidroméis mais fortes podem apresentar uma leve nota alcoólica picante, e os carbonatados tendem a expressar mais aroma. Os aromas provenientes da fruta devem complementar e realçar o hidromel base, não sobrepujá-lo. Algumas frutas (e seus caroços) podem adicionar um aroma semelhante ao de especiarias.

\textbf{Aparência}: Limpidez de boa a brilhante; quanto maior a limpidez, mais desejável. A presença de bolhas, espuma, colarinho ou efervescência depende do nível de carbonatação declarado. A cor é derivada da variedade de fruta e de mel utilizadas. Quanto maior a vivacidade e o brilho da cor, mais desejável.

\textbf{Sabor}: Similar ao Aroma em termos de caráter e intensidade de fruta e mel, tornando-se mais evidente nas versões mais fortes e doces. O caráter da fruta deve ser perceptível e estar em equilíbrio com o nível de dulçor do hidromel. Algumas frutas podem trazer sabores básicos (amargo, doce, azedo), além de suas características próprias. Muitas frutas acrescentam acidez e taninos naturais, que precisam estar em equilíbrio com o hidromel base.

O dulçor residual e o final de boca devem corresponder ao nível de dulçor declarado. A acidez deve ser equilibrada, suave ou vibrante, mas não agressiva. Os taninos podem fazer um hidromel mais doce parecer mais seco. Características de fermentação ásperas, desagradáveis, excessivamente sulfurosas ou de levedura são indesejáveis.

\textbf{Sensação na Boca}: A carbonatação segue o nível declarado, podendo ir de tranquilo a espumante. O corpo tende a aumentar com o dulçor e a força alcoólica, variando geralmente de médio-baixo a alto. A percepção alcoólica acompanha a força declarada, podendo ser desde ausente até perceptível e com sensação de aquecimento. Os taninos de frutas mais escuras podem conferir algum corpo e uma sensação adstringente, que não deve ser extrema.

\textbf{Ingredientes}: Os mesmos ingredientes de um M1 Traditional Mead, com a adição de pelo menos um tipo de fruta de caroço.

\textbf{Comentários}: Um Stone Fruit Mead que também contenha outra fruta deve ser inscrito como M2E Other Fruit Mead. Versões com especiarias devem ser inscritas como M3C Fruit and Spice Mead. Hidroméis com frutas e com adição de outros ingredientes que não sejam frutas provavelmente devem ser inscritos como M4F Experimental Mead.

\textbf{Instruções para Inscrição}: Os participantes devem especificar os níveis de dulçor, carbonatação e força alcoólica. As variedades de mel podem ser declaradas. A fruta utilizada deve ser declarada.

\textbf{Exemplos Comerciais}: Moonlight Entice, Prairie Rose Chokecherry, Redstone Peach Reserve, Redstone Sunshine Nectar, Schramm's Mead The Statement, White Winter Honey of a Plum.
